% !PRINTED 472
\weberpage{472}

%<<CONT
one sets $h_x = 0$. The quantities $u_x, g_x, h_x$ depend only on the
variables $x_1, x_2, x_3$.

Now, as \gesperrt{Klein} has remarked loc.\,cit.\ ${}^1$), the solution
of the general resolvent can be reduced to that of the special one,
if one adjoins the root of a \gesperrt{biquadratic equation}.

To prove this, we introduce beside the point $(x)$
a second point $(y)$ and form the polar of $f$:
\weq{3}{f_1(x,y) = y_1f'(x_1) + y_2f'(x_2) + y_3f'(x_3).}
If we then subject $(x)$ and $(y)$ simultaneously to the same substitution
of the group $G_{168}$, the function $f_1(x,y)$ remains
unchanged (Bd.\ I, \S.\,60). We now take the point $(x)$ arbitrarily,
but require of the point $(y)$ that it shall lie simultaneously
on the fundamental curve and on the polar of the point $(x)$,
so that simultaneously
\weq{4}{f(y_1,y_2,y_3) = 0, \quad f_1(x,y) = 0}
shall hold. Then to each point $x$ there correspond four points $y$, and
if we put for $(x)$ a point connected with it, then
each of these four points $(y)$ likewise passes into a connected
point. The function $h_y$ is now $=0$ and $g_y$ is a four-valued
function of the point $(x)$. Symmetric functions
of these four values remain unchanged under the substitutions
of $G_{168}$, and consequently $g_y$ is the \gesperrt{root of a biquadratic
equation whose coefficients depend rationally on the $g_x, h_x$}.
The root of this biquadratic equation must be
adjoined.

The function $u_y$ is a root of the special resolvent
\weq{5}{R(u_y, g_y, 0) = R_y = 0.}
We now denote the four points $y$ belonging to the same $x$
by $y, y', y'', y'''$ and form the symmetric function of these
four points
\weq{6}{(t-u_y)(t-u_{y'})(t-u_{y''})(t-u_{y'''}) = \Phi(t)}
for an indeterminate $t$. This function remains unchanged under
all substitutions of the group $G_{24}$, and is therefore rationally
expressible in terms of $u_x, g_x, h_x$.

\par\noindent\rule{0.3\linewidth}{0.4pt}\par
${}^1$) Mathematische Annalen, Bd.\ XV, S.\ 280.
